% Korean edition of sections/s4_1_tests.tex.
\seclabel
\claim{\secnum{4.1}{sec:wrappertests}\quad 시작, 피드, 매수, 매도는 동작했고, \code{plan()}에는 100\% 상한이 필요했다}

\lead{\ref{sec:design}장의 설계로 만든 래퍼를 2026년 11월 24일 09:27\textasciitilde 09:55 KST에 실서버에서 실행했고, \ref{sec:oneshare}절의 내림 테스트는 10:26과 10:33에 했다.} 패키지 \code{pytimefolio}에는 \figref{modules}의 모듈이 들어 있다. 전략은 \code{client} 모듈의 클래스 \code{Timefolio} 하나를 쓴다(\figref{client}). 가격 인자(\code{**price})는 \code{limit_prc}, \code{limit_idx}, \code{stop_prc} 중 정확히 하나이며, \code{hm0}과 \code{hm1}은 선택이다(\figref{orderapi}).

% Korean edition of figs/client.tex.
\def\fclL#1#2{\hspace*{#1mm}\code{#2}}% one code line, indented #1 mm
\def\fclC#1{\leavevmode{\rmfamily\color{greyD}#1}}% comment beside a code line
\fig[wide]{client}{클래스 하나가 전략에 상태, 주문, 가이드라인 현황을 준다}[\code{Timefolio} 클라이언트 0.1.0 사용법, 오른쪽은 설명]{%
{\fontsize{7.6}{9.6}\selectfont\begin{tabular}{@{}>{\RaggedRight\arraybackslash}p{106mm} >{\RaggedRight\arraybackslash}p{70mm}@{}}
\fclL{0}{async with Timefolio(91, email=os.environ["EMAIL"],} & \fclC{대회 91. 비밀번호는 넘기기만 하고}\\
\fclL{28.4}{password=os.environ["PASSWORD"]) as tf:} & \fclC{저장하지 않음. 쿠키 묶음은 저장됨}\\[2pt]
\fclL{4}{w = tf.weights()} & \fclC{\code{prodId} $\rightarrow$ 현재 \code{wei}(\%)}\\
\fclL{4}{tf.orders_left()} & \fclC{80에서 오늘의 새 주문과 정지를 뺀 값}\\[2pt]
\fclL{4}{orders = tf.plan(targets, min_wei=0.05, max_total=100,} & \fclC{dry run: 아무것도 보내지 않음}\\
\fclL{27}{prices={"A000660": p}, limit_idx=5)} & \fclC{규칙 위반이면 \code{ValueError}}\\
\fclL{4}{tf.dropped} & \fclC{1주 미만이라 뺀 주문}\\[2pt]
\fclL{4}{try:} & \\
\fclL{8}{ids = await tf.rebalance(targets, limit_idx=5)} & \fclC{\code{plan()} 후 1\,s 간격으로 주문}\\
\fclL{4}{except RebalanceError as e:} & \\
\fclL{8}{e.placed} & \fclC{실패 전에 보낸 주문 ID}\\[2pt]
\fclL{4}{i = await tf.buy("A005930", 0.1, limit_prc=150700)} & \fclC{또는 \code{tf.sell}. 주문 ID를 반환}\\
\fclL{4}{await tf.stop(i)} & \fclC{\code{stop(*ids)}는 여러 개를 받음}\\
\fclL{4}{r = await tf.wait(i)} & \fclC{\code{r.outcome}, \code{r.cum_qty}, \code{r.qty2ord}, \code{r.avg_px}}\\[2pt]
\fclL{4}{g = await tf.guidelines()} & \fclC{\code{g["ok"]}, \code{g["sector"]}, \code{g["turnover"]}}\\
\end{tabular}}}{\code{targets}는 종목코드에서 목표 비중(\%)으로 가는 딕셔너리이며 포트폴리오 전체를 나타낸다. 목표에서 빠진 보유 종목은 목표 0이 되어 매도된다. \code{plan()}은 진행 중인 주문이 아직 채울 몫을 뺀 뒤 \code{rebalance()}가 보낼 주문을 매도 먼저 돌려준다. \code{prices}는 포트폴리오가 보유하지 않은 종목의 가격으로, 1주 검사를 할 수 있게 한다(\ref{sec:oneshare}절). 가격 인자(\code{**price})는 \code{limit_prc}, \code{limit_idx}, \code{stop_prc} 중 정확히 하나이며, \code{hm0}과 \code{hm1}은 선택이다. \code{AddOrder}가 주문 ID를 돌려주지 않으므로 \code{buy()}와 \code{sell()}은 피드에서 찾은 주문 ID를 돌려준다(\ref{sec:orders}절). \code{rebalance()}는 먼저 서버에서 오늘 주문을 다시 세고, 주문 수가 하루 한도의 남은 건수를 넘으면 아무것도 보내지 않는다. \code{r.outcome}은 \code{filled}, \code{stopped}, \code{partly_filled_then_stopped} 중 하나다. \code{p}는 호출자가 주는 원화 가격이다.}{DART Quant Team; pytimefolio 0.1.0의 \code{client.py}}


\lead{\code{async with} 블록에 들어가면 전략이 제어를 받기 전에 5단계가 실행된다.} 첫째, 클라이언트는 \code{session.json}에 저장된 쿠키 묶음으로 액세스 토큰을 갱신하고, 갱신이 실패할 때만 로그인한다(D17). 둘째, \ref{sec:queries}절의 규칙으로 포트폴리오를 고른다. 셋째, 실시간 피드를 백그라운드 태스크로 시작하고 첫 스냅샷을 최대 30\,s 기다린다. 오지 않으면 타임아웃으로 시작에 실패하므로 \code{weights()}가 빈 상태를 읽는 일은 없다. 이후 상태를 읽거나 주문을 넣는 모든 호출은 마지막 피드 응답이 26\,s(2 $\times$ 3\,s + 20\,s)보다 오래되면 \code{StaleFeed}를 던지므로, 낡은 비중을 쓰는 일도 없다. 넷째, \code{Universe}를 한 번 불러와 종목별 섹터를 익힌다(\ref{sec:guidelines}절). 다섯째, \code{Portfolio/Orders}에서 그날의 주문(새 주문과 정지)을 세고, 그 수와 자체 건수 파일 중 큰 값으로 하루 한도를 시작한다. 그래서 하루 중에 다시 시작해도 건수가 0으로 돌아가지 않는다. 블록을 나가면 피드 태스크를 취소하고 WebSocket과 HTTP 클라이언트를 닫으며, \code{Logout}은 호출하지 않는다(D16).

\Needspace{3\baselineskip}\lead{각 모듈이 실서버 테스트를 통과했고, 오프라인 테스트 31개가 계산과 실서버에서 일으킬 수 없는 실패 경로를 다룬다(\figref{tests}).} 오프라인 테스트 중 9개는 순수 함수를 검사한다. 나머지 22개는 서버를 \code{httpx.MockTransport}와 가짜 객체로 대체해, 실서버가 원할 때 만들어 주지 않는 상황을 만든다. 오류 상태 코드별 응답, \code{AddOrder} 응답 유실, 재전송도 실패하는 경우, 핑만 보내는 피드, 낡은 피드, 동시 쓰기 3건(모의 서버에서 1.002\,s 간격), 연속된 같은 주문 2건, 주문 하나에 맞는 새 레코드 2개, 그리고 읽기나 쓰기에서 403이나 429를 받은 뒤 200이 오는 경우다. \figref{tests}의 주에 전체 목록이 있다. 31개 모두 push와 pull request마다 GitHub Actions에서 Python 3.11\textasciitilde 3.13으로 실행된다. \code{tests/live}의 실서버 검사는 CI에서 실행하지 않는다.

% Korean edition of figs/tests.tex.
\fig[wide]{tests}{실서버 테스트에서 드러난 결함 2개를 고친 뒤 모든 테스트가 통과했다}[목표(\figref{goals})별 래퍼 테스트: 2026년 11월 24일 09:27\textasciitilde 10:34 KST 대회 서버 대상 실서버 테스트(포트폴리오 하나)와 오프라인 테스트]{%
{\tabfont\begin{tabularx}{\linewidth}{>{\RaggedRight\arraybackslash}p{16mm} >{\RaggedRight\arraybackslash}p{25mm} >{\RaggedRight\arraybackslash}p{47mm} >{\RaggedRight\arraybackslash}X}
\thd{목표} & \thd{테스트} & \thd{검사한 것} & \thd{결과}\\
G1 & \code{transport} & 첫 요청, 두 번째 클라이언트의 첫 요청, 재사용 연결의 비용. \code{Date} 헤더로 구한 시계 오차 & 새 프로세스의 첫 요청 1.2\textasciitilde 1.3\,s. 두 번째 클라이언트의 첫 요청 130\ms. 재사용 31\textasciitilde 39\ms. 서버 시계 빼기 PC 시계 $-2.98$\,s로 \ref{sec:auth}절의 3\,s와 일치\\\rowsep
G1 & \code{auth} & 두 번 실행하며 로그인, 쿠키 저장, 갱신 & 첫 실행은 로그인하고 쿠키 묶음을 \code{~/.config/pytimefolio/session.json}에 권한 0600으로 저장. 두 번째 실행은 로그인 없이 토큰 갱신\\\rowsep
G2, G5 & \code{rest} & 포트폴리오 규칙. 연결 하나에서 동시 조회. 종목코드 검사 & 규칙으로 운용 중인 포트폴리오를 선택. \code{Session}, \code{Universe}(2,648), \code{Trades}(64), \code{Orders}(그날 오전 5건)를 HTTP/2 연결 하나에서 동시에 가져옴. \code{A005930}의 \code{ValidateProduct}는 경고 없음\\\rowsep
G2 & \code{hub}, \code{state} & 3\,s마다 폴링한 연결 하나. 델타 병합 & 스냅샷 1개 뒤 델타 20개, 끊김 없음. 병합 상태가 \code{t}가 같은 스냅샷과 15회 중 15회 일치(\ref{sec:feed}절의 테스트)\\\rowsep
G3 & \code{orders}, 매수 & 삼성전자 0.1\%, 지정가 150,700원, 시장가 188,400원보다 20\% 낮음 & \code{AddOrder} 후 1\textasciitilde 3\,s 만에 피드에서 찾음: Id 1650412, \code{Working}, \code{qty2ord} 4. 3\,s 뒤 정지. \code{wait()}은 정지, 4주 중 0주 반환\\\rowsep
G3 & \code{orders}, 매도 & LIG넥스원 0.1\%, 지정가 427,000원, 시장가 355,500원보다 20\% 높음 & 주문 레코드 \code{ls} = L, \code{ex} = X, \code{sgn} = $-1$, \code{qty2ord} 2. 정지. 결과 정지, 2주 중 0주\\\rowsep
\hl G3 & \code{plan()}, 합계 & 삼성전자 50\%, \code{weiLimit} 40\% 대비 & 40\%로 상한 적용. water-filling이 초과분을 다른 보유 종목 22개에 분산. 목표 합계 139\%를 \code{plan()}이 받아들임: 결함 1\\\rowsep
G3 & \code{plan()}, 1주(\ref{sec:oneshare}) & \code{min_wei} 0.01로 작은 주문 3건, 09:52 & SK하이닉스 +0.03\%는 경고와 함께 뺌(1주 0.1413\%). 삼성전자 +0.1\%와 리노공업 +0.06\%는 유지\\\rowsep
\hl G3 & \code{guidelines()}(\ref{sec:guidelines}) & 피드 필드와 \code{Violations} 응답, 09:52 & 서버는 위반 0건인데 IT가 섹터 한도를 넘은 것으로 표시: 결함 2. 고친 뒤 서버 위반 0건. 회전율 4.8\%로 \code{ok} False\\\rowsep
G3, G5 & \code{plan()}, 섹터(\ref{sec:guidelines}) & 목표 3가지, 09:55 & 비중 그대로: 주문 0건. 산업재 +6\%p: 거부. 산업재 +3\%p: 주문 2건(\figref{sectortest})\\\rowsep
G2 & \code{live_fields} & 10:33에 60\,s 동안 델타마다 \code{bs_updated}. 주문 없음 & 델타 20개 중 20개에 \code{bs_updated}가 딕셔너리로 옴. \code{na}가 844,505,167원에서 844,076,987원으로 움직여 순자산이 최신으로 유지됨\\\rowsep
G3 & \code{live_rounding}(\ref{sec:oneshare}) & 삼성전자 0.04\%(1.76주어치), 시장가보다 20\% 낮은 지정가. 즉시 정지 & 두 번 모두 \code{qty2ord} 1(Id 1650693 약 10:26, Id 1650719 10:33:45), 0주 체결: 서버는 내림함\\\rowsep
G3 & 수정된 코드, 실서버 & 리뷰 뒤 바꾼 코드로 주문을 넣고 정지: 크기를 0.01\%로 맞춤. 서버가 받았을 수 있는 주문은 재전송 안 함 & 주문 레코드에 \code{wei}가 보낸 그대로 저장(0.04, 0.1, 0.3 확인). 주문 1650719는 피드에서 \code{Accepted}로 찾았고, 즉시 정지했으며, \code{Done}으로 끝남. 10:34 카운터: 새 주문 9건과 정지 7건, 80건 중 16건 사용\\\rowsep
G1\textasciitilde G3, G5 & 오프라인 & 단위 테스트 31개(순수 9개, 모의 HTTP로 실패 경로 22개) & 모두 통과\\
\end{tabularx}}}{시장가는 주문 시점 피드의 가격이다. 테스트 주문은 체결되지 않도록 시장가에서 20\% 떨어진 지정가였다. 강조 행이 결함 2개다. 결함 1: water-filling은 합계를 유지하므로 139\% 목표가 상한 적용 뒤에도 139\%였다. 이제 \code{plan()}은 목표 합계가 \code{max_total}(기본 100\%)을 넘으면 \code{ValueError}를 던진다. 결함 2: 섹터 한도를 섹터 코드 \code{sec}가 아니라 \code{msecs} 레코드 \code{Id}로 찾아 IT를 다른 섹터의 한도와 비교했다. \code{sec}로 찾게 고쳤다. 09:52의 \ref{sec:oneshare}절 테스트와 \ref{sec:guidelines}절 테스트는 \code{plan()}과 \code{guidelines()}만 호출했고 주문은 보내지 않았다. 오프라인 순수 테스트: 상태 병합(문자열 키를 \code{int}로 변환, 바뀌지 않은 필드 유지, 널 응답, 추가·삭제 레코드, 스냅샷 교체), water-filling, 목표-주문 변환, 09:18에 캡처한 요청과 주문 본문 비교, \code{sector_violations}, KRX 호가 단위. 실패 경로(\code{httpx.MockTransport}와 가짜 객체): 상태 코드 401, 403, 429, 400, 502를 예외로 변환. 전송 계층은 연결 실패 때 GET만 한 번 재시도하고 POST는 재시도하지 않으며, \code{AddOrder}는 연결이 맺어지지 않았을 때만 한 번 더 보냄. 읽기 오류는 보낸 것으로 셈. 동시 \code{token()} 호출 5건이 \code{SilentRefresh} 1건. 401이면 한 번 갱신. 손상된 세션 파일이면 로그인. 크기 맞춤. 진행 중 주문은 다시 주문하지 않음. 하루 한도와 동기화. 연결 실패는 세지 않음. 동시 쓰기 3건이 최소 1\,s 간격으로 서버에 도달(1.002\,s 측정). 403과 429 각각이 모든 쓰기를 멈추며, 서버가 다시 200을 줘도 마찬가지. \code{AddOrder} 응답 유실 시 재전송 안 함. 낡은 피드는 주문과 \code{wait()}을 막음. 피드 신선도. 실서버 테스트 뒤 추가: 재시작해도 쿠키 만료 유지. 재전송이 실패해도 검색하며 세 번째 전송 없음. 같은 주문을 두 번 넣으면 ID 2개. 주문 하나에 맞는 새 레코드 2개면 \code{AmbiguousState}. 읽기에서 403이면 쓰기 중단. 서버 핑은 응답 마감을 늘리지 않음. 차단된 피드는 재연결을 멈춤. 이 테스트들은 CI에서 Python 3.11\textasciitilde 3.13으로 돌며, \code{tests/live}는 돌지 않는다.}{DART Quant Team; pytimefolio 0.1.0의 \code{tests/}와 \code{tests/live/}}


\lead{이 테스트 전 코드 리뷰에서 주문과 피드 경로의 결함 7개를 찾았고, 테스트를 돌린 코드에서는 7개 모두 고쳐져 있었다.} 고친 내용은 이미 설명한 동작들이다. (1) \code{wei}를 0.01\%로 한 번 반올림하고 보낸 값으로 대조, (2) 응답을 잃은 주문을 자동으로 다시 보내지 않고 쓰기와 검색 내내 잠금 하나 유지(\ref{sec:orders}절), (3) 26\,s 뒤 \code{StaleFeed}, (4) full-jitter 백오프로 재연결, (5) \code{bs_updated} 병합(\ref{sec:feed}절), (6) \code{plan()}에서 진행 중인 주문 빼기(\ref{sec:guidelines}절), (7) 정지를 포함하고 재시작에도 유지되는 하루 건수, 그리고 403이나 429 뒤 모든 쓰기 중단(\ref{sec:design}장)이다.

\lead{매도 테스트로 매도를 보내는 방식이 정해졌다.} 래퍼의 \code{sell()}은 매수와 같은 \code{AddOrder} 본문에서 청산 코드만 E 대신 \code{ex} = X로 보내며, \code{ls}는 L 그대로다. 서버는 이를 받아들여 부호 \code{sgn} = −1(매도), \code{qty2ord} 2로 기록했다. 따라서 매도는 매수와 같은 방향 코드(\code{ls} = L)에 청산(\code{ex} = X)으로 보낸다. \ref{sec:orders}절의 주문은 모두 매수였으므로, 이 테스트로 양방향 모두 넣고, 피드에서 찾고, 정지해 본 셈이다.

\lead{첫 \code{plan()} 테스트에서 합계가 100\%를 넘는 목표가 그대로 통과함을 발견했다.} water-filling은 한도에서 잘린 비중을 다른 종목으로 옮기고 합계는 유지하므로, 139\% 목표는 상한 적용 뒤에도 139\%였고 \code{plan()}은 자산의 139\%에 해당하는 포트폴리오를 요구할 뻔했다. 이제 \code{plan()}은 상한을 적용하기 전에 목표 합계가 \code{max_total}(기본 100\%)을 넘으면 \code{ValueError}를 던진다(D19).

\lead{10:34 기준 그날 건수는 새 주문 9건과 정지 7건이었고, 래퍼 카운터는 80건 중 16건을 썼다.} \ref{sec:orders}절의 테스트가 주문 5건과 정지 3건을 썼고, 매수와 매도가 각각 2건씩, \ref{sec:oneshare}절의 내림 테스트가 각각 2건씩 더했다. \code{plan()}과 \code{guidelines()} 테스트는 아무것도 보내지 않았다. 취소가 100건에 포함되는지 주최측이 답할 때까지(\figref{openq}의 질문 Q2, \ref{sec:validity}절) 래퍼는 정지를 주문으로 센다. 따라서 래퍼 한도 아래로 64건이 남고, 주최측 기준 100건까지는 84건이 남는다.

\lead{테스트는 G2와 G5를 보여주며, G1과 G3는 일부만, G4는 아직 확인하지 못했다(\figref{goalstatus}).} G1은 래퍼가 11월 30일까지 로그인 상태를 유지해야 하는데 테스트는 오전 한나절이었다. 두 번째 실행은 로그인 없이 저장된 쿠키 묶음으로 갱신했지만, 한 실행 안에서 720\,s를 넘겨 갱신한 프로세스나 7일 동안 돈 프로세스는 아직 없다. G2는 두 부분으로 판단한다. 병합 상태가 \code{t}가 같은 스냅샷과 15회 중 15회 일치한 것은 실서버에서 관찰했다(\ref{sec:feed}절). 26\,s보다 오래된 상태를 거부하는 것은 실제 장애를 일으키지 않았으므로 오프라인에서만 보였다. G3는 일부만 보였다. 15\,s 확인 시간은 1\textasciitilde 3\,s 안에 찾은 실제 매수에 근거하며 느린 서버에서는 시험하지 않았고, \code{plan()}은 실행 중의 비중이 아니라 계획 시점의 목표 비중을 검사하기 때문이다. G4, 즉 폴링이 주문을 늦추지 않는다는 목표는 아직 결과 측정 기준이 없다. 테스트는 피드와 주문이 이벤트 루프 하나를 공유한다는 것만 보여준다. G5는 실서버·오프라인 테스트를 모두 통과했다.

% Korean edition of figs/goalstatus.tex.
\fig[wide]{goalstatus}{테스트는 G2와 G5를 보여주며, G1과 G3는 일부만, G4는 측정하지 않았다}[\figref{goals}의 목표마다 관련된 \figref{tests}의 테스트, 2026년 11월 24일. G2의 두 부분은 따로 판단]{%
{\tabfont\begin{tabularx}{\linewidth}{>{\RaggedRight\arraybackslash}p{34mm} >{\RaggedRight\arraybackslash}p{46mm} >{\RaggedRight\arraybackslash}X >{\RaggedRight\arraybackslash}p{15mm}}
\thd{목표} & \thd{테스트} & \thd{결과} & \thd{상태}\\
\hl G1 사람 없이 대회 기간 내내 로그인 유지 & 프로세스 2개에 걸친 \code{auth} 실서버 테스트. 오프라인: 동시 \code{token()} 호출 5건, 401 재시도, 손상된 세션 파일, 재시작 뒤 쿠키 만료 & 두 번째 프로세스가 로그인 없이 저장된 쿠키 묶음으로 갱신. 오프라인에서 호출 5건이 \code{SilentRefresh} 1건, 만료가 재시작 뒤에도 유지. 한 프로세스 안에서 720\,s 뒤 갱신과 7일 쿠키 만료 전 로그인은 아직 실행해 보지 않음 & 일부 확인\\\rowsep
G2(a) 병합: 병합 상태가 \code{t}가 같은 스냅샷과 같음 & \code{hub}, \code{state}, \code{live_fields} 실서버 테스트 & 병합한 \code{pos}와 \code{ords}가 \code{t}가 같은 스냅샷과 15회 중 15회 일치. 델타 20개 중 20개의 \code{bs_updated}가 \code{na}를 최신으로 유지. \code{bs}와 \code{msecs}는 필드 단위로 비교하지 않음 & 확인\\\rowsep
G2(b) 신선도: 26\,s보다 오래된 상태로 읽기나 주문을 하지 않음 & 오프라인: 낡은 피드 테스트, 피드 신선도, 응답 마감을 늘리지 않는 핑 & 마지막 응답이 26\,s보다 오래된 피드는 \code{StaleFeed}를 던지고 주문과 \code{wait()}을 막음. 핑만 보내는 피드는 응답 마감 뒤 타임아웃. 실제 장애는 일으키지 않음 & 오프라인 확인\\\rowsep
\hl G3 주최측 한도와 대회 가이드라인 안에서 주문을 넣고, 정지하고, 추적 & 오프라인: 하루 한도와 동기화, 쓰기 간격, 재전송 안 함, 재전송 실패, 같은 주문 2건, 403·429에서 중단. 실서버: 매수, 매도, \code{plan()}, 수정된 코드 & 한도를 넘는 쓰기 거부. 잃어버린 \code{AddOrder} 응답은 재전송 안 함. 동시 쓰기 3건이 모의 서버에 1.002\,s 간격으로 도달. 403이나 429 뒤 모든 쓰기 중단. 실제 주문을 넣고 1\textasciitilde 3\,s 만에 찾아 정지, 10:34까지 80건 중 16건 사용. 보이지 않은 것: 느린 서버에서 15\,s 안 확인(실제로 찾은 것은 매수뿐), 실행 중 한도(\code{plan()}은 계획 시점의 목표 비중만 검사) & 일부 확인\\\rowsep
\hl G4 폴링이 주문을 늦추지 않음 & 없음 & 피드 폴링과 주문이 이벤트 루프 하나를 공유. 폴링이 진행 중일 때 주문의 지연은 측정하지 않음 & 미확인\\\rowsep
G5 모호하면 추측하지 않고 멈춤 & \code{rest} 실서버(포트폴리오 규칙). 오프라인 응답 유실, 재전송 실패, 상태 코드 테스트. \code{plan()} 섹터·합계 테스트 & 규칙으로 포트폴리오 선택. 찾지 못한 주문은 재전송 없이 \code{AmbiguousState}. 400은 \code{RequestError}. \code{plan()}이 산업재 한도를 넘는 목표와, 결함 1 수정 뒤 100\% 초과 목표를 거부 & 확인\\
\end{tabularx}}}{강조 행은 완전히 확인되지 않은 목표다. 확인: 나열한 테스트가 \figref{goals}의 성공 기준을 충족함. 오프라인 확인: 모의 서버에서 코드 동작을 보였지만 실서버에서는 관찰하지 않음. 일부 확인: 기준의 일부를 충족했고 나머지는 아직 실행해 보지 않음. 미확인: 목표를 측정한 테스트가 없음. 각 결과로 사용자가 무엇을 어떤 가정 아래 믿을 수 있는지는 \figref{guarantees}에 있다. 실서버 테스트는 계정 하나로 오전 한나절 동안 했다(\ref{sec:validity}절).}{DART Quant Team; \figref{tests}}


\verdict{래퍼의 \code{plan()}은 합계가 100\%를 넘는 목표를 거부한다(D19). \code{sell()}은 실서버에서 확인한 대로 \code{ls} = L, \code{ex} = X로 보낸다(D21).}
